\section{Heaven and Earth Are Silent}

\paragraph{Sonnet 164 by Petrarch}


This is my more esoteric translation, in iambic pentameter, of Sonnet 164 by Petrarch.

\begin{verse}
Hor che'l ciel e la terra e'l vento tace,\\
e le fere e gli augelli il sonno affrena,\\
notte il carro stellato in giro mena,\\
e nel suo letto il mar senz'onda giace.

Veglio, penso, ardo, piango; e chi mi sface\\
sempre m'è innanzi per mia dolce pena;\\
guerra è'l mio stato, d'ira et di duol piena,\\
e sol di lei pensando ho qualche pace.

Così sol d'una chiara fonte viva\\
move'l dolce e l'amaro ond'io mi pasco;\\
una man sola mi risana e punge.

E perché'l mio martir non giunga a riva,\\
mille volte il dì moro e mille nasco,\\
tanto dalla salute mia son lunge.
\end{verse}

\begin{verse}
Now that Heaven, Earth and Wind are silent\\
The beasts and the birds are subdued by sleep,\\
Night leads the starry chariot away,\\
The waveless sea lies sleeping in its bed,

I awake, think, burn, weep: who destroys me\\
Stands before my eyes to my sweet sorrow:\\
War is my state, filled with grief and anger,\\
Thinking only of her brings me some peace.

Only from a pure and living fountain\\
Flow the sweet and bitter on which I feed;\\
The same hand both heals me and pierces me.

And because my torment cannot reach shore,\\
I am born and I die a thousand times,\\
I remain so far from my salvation.
\end{verse}

\begin{quotex}
Heaven, Earth, and the Wind represent the Three Worlds.

The Sea is the soul, which needs to become calm and free of perturbations.

The animals are the warring parts of his inner life.

It was set to music by Claudio Monteverdi.
\end{quotex}

\url{https://www.youtube.com/watch?v=mPhWDKT2gSk} 

\flrightit{Posted on 2021-10-15 by Cologero}

